\section{Ontology and Vocabulary Catalogue}
\label{sec:appendix-ontology}

\subsection{Where the ontology lives}

There is no single ontology file. Authority is composed across four modules, each
owning a vocabulary, unioned at read time: a frozen semantic seal (node types and
the closed predicate list), an enrichment layer (textual and structural predicates,
plus the trust envelope), a domain layer (labels, predicates, endpoint signatures),
and a grading layer (the per-edge grade and the attribution contract). A fifth
vocabulary --- the reader-facing one exposed by the \mbox{API} --- is deliberately
distinct from the writer's (\cref{sec:collision}).

The composition function carries the method's key statement, which we quote because
it generalises beyond this project:

\begin{quote}\small
Critically, the graph is not an input. The previous scorecard gate computed
\code{declared = ontology | everything in the graph} and then asked which graph types
were missing from \code{declared}; nothing can be, so it reported 0 undeclared for a
whole wave while eleven predicates went unclassified. \textbf{The data may not
declare itself valid.}
\end{quote}

\subsection{The four typing axes}

\begin{description}[leftmargin=1.6em,itemsep=2pt,topsep=3pt,style=nextline]
\item[\code{KnowledgeLayer}] \code{L1\_SOURCE\_EXPLICIT}, \code{L2\_DETERMINISTIC\_DERIVED},
\code{L3\_LLM\_EXTRACTED}, \code{L4\_INTERPRETIVE\_CLAIM}. Documented as ``never
collapsed, never inferred from a score''.
\item[\code{QualityTier}] \code{TIER\_A} to \code{TIER\_D}, a total bijection from the
layer. One further transition: an \textsc{l3} claim in state \code{CANDIDATE} is
\code{TIER\_D}, not \code{TIER\_C}, because \code{TIER\_C} means model-extracted
evidence that survived review.
\item[\code{AttributionPrecision}] \code{PER\_PASSAGE},
\code{CONTAINER\_INHERITED}, \code{TEXTUAL\_MENTION}, \code{NOT\_AN\_ATTRIBUTION}.
The fourth member is what makes a graph-wide universal-quantifier invariant over this
property expressible at all.
\item[Evidence surface] \code{SANSKRIT}, \code{STRUCTURAL}, \code{SOURCE\_METADATA},
\code{MIXED}, \code{TRANSLATION}, \code{SHARED\_REGISTRY\_ENTITIES}. Stored under
the property name \code{evidence\_basis}, which collides with an unrelated derivation
enumeration of the same name (\cref{sec:collision}).
\end{description}

Two legacy spellings of the layer field survive from an earlier pipeline generation
(\code{SOURCE\_EXPLICIT} and \code{DETERMINISTIC\_DERIVED} without the
\code{L}-prefix), on 6{,}236 and 8{,}214 edges respectively. The fact freeze reports
both raw and folded distributions.

\subsection{Lifecycle vocabularies}

Six lifecycle vocabularies exist, deliberately not unified, each scoped to one layer.
The two that matter for this paper are the graph-edge state
(\code{CANDIDATE}/\code{ACCEPTED}/\code{REJECTED}, with the promotion gate encoded as
a frozen set of trust classes) and the semantic-artifact status
(\code{CANDIDATE}, \code{AUTO\_ACCEPTED}, \code{HUMAN\_ACCEPTED},
\code{HUMAN\_REJECTED}, \code{VALIDATION\_REJECTED}, \code{SUPERSEDED},
\code{NEEDS\_REVIEW}), of which a model may write only the first.

Rejections carry one of twenty codes, one code per failure mode, including
\code{EVIDENCE\_TOKEN\_NOT\_IN\_PASSAGE}, \code{EVIDENCE\_HASH\_MISMATCH},
\code{PREDICATE\_FORBIDDEN}, \code{PREDICATE\_NOT\_WHITELISTED},
\code{EXPLICITNESS\_UNSUPPORTED\_BY\_EVIDENCE} and \code{DUPLICATE\_ASSERTION}.

Review verdicts applied to candidates are \code{ACCEPT\_MODEL\_REVIEWED} (promoted to
\code{TIER\_C}, \code{review\_state} set to \code{MODEL\_ADJUDICATED} and never to
\code{HUMAN\_REVIEWED}), \code{REJECT} (the edge is deleted, since a reader running
the obvious query never sees the tier, and the row stays on disk so the deletion is
reproducible), and \code{AMBIGUOUS} / \code{NEEDS\_MORE\_EVIDENCE} (kept at
\code{TIER\_D} and stamped).

There is no \code{WITHDRAWN} state. \code{SUPERSEDED} is the nearest.

\subsection{Withholding mechanisms}

\code{WITHHELD} is not a member of any enumeration. It is four separate typed
mechanisms, each scoped to a layer: a reason string per unmapped sealed predicate,
emitted into the load report alongside a count of predicates that are neither mapped
nor explicitly withheld; a typed node property with a declared class space
(S\=amavedic notation, three classes); typed non-resolution rows in the lexical layer
(711 tokens, three linguistic reasons); and edge-type-keyed withholding for
near-parallel pairs.

\subsection{Node labels and relationship types}

55 labels are in use and 81 relationship types, of a token table declaring 56 labels
--- \code{SourceArtifact} is declared and carries no nodes. The full distributions
are in \code{supplementary\vgslash census\_node\_labels.csv} and
\code{supplementary\vgslash census\_relationship\_types.csv}; the fourteen largest of each
are in \cref{tab:labels}.

\subsection{Grades the bulk stamper must not touch}

Twenty-one predicates are excluded from the generic grade sweep, each with a recorded
reason. Three classes of reason recur. A predicate whose grade is the output of an
independent review cannot be reproduced by any provenance signature --- a post-load
regrade once silently erased 587 promotions in one pass. A predicate whose grade is
the grade of the node across the edge cannot be computed by a stamper that cannot see
across an edge. And a predicate whose tier depends on which derivation produced the
row must have both directions of its mirror written byte-identically, or a query gets
a different tier depending on which way it walked.
